% =============================================================================
% 02_RELATED_WORK — Thematic Literature Review
% =============================================================================
% Page target: ~2 pages
% Structure: Thematic organization — do NOT copy abstracts
% Key instruction: synthesize, don't catalogue. Show analytical narrative.
% =============================================================================

\section{Related Work}
\label{sec:related_work}
\addcontentsline{toc}{section}{Related Work}

The existing literature relevant to this review can be organized into three
interlocking thematic streams: (\textit{i}) technical approaches to AI
evaluation and validation; (\textit{ii}) regulatory and governance frameworks
for clinical AI; and (\textit{iii}) implementation science perspectives on
digital health and AI integration. These streams are not independent — the
evaluation literature is shaped by regulatory definitions of valid performance,
and the governance literature is increasingly informed by implementation
outcomes in real clinical settings. This section draws out the analytical
connections across these streams.

% =============================================================================
% 2.1 AI EVALUATION — technical validation methods and their limits
% =============================================================================

\subsection{AI Evaluation and Validation Methods}
\label{sec:ai_evaluation}

The technical evaluation of health AI has grown into a mature subfield,
producing consensus guidelines on benchmarking, external validation, and
reporting standards \citep{pham2025ethical}. However, a consistent weakness
across the evaluation literature is the reliance on performance metrics
computed in silico on curated retrospective datasets — an approach that
systematically overstates real-world reliability. \citet{abulibdeh2025illusion}
document this gap sharply in their critical analysis of FDA-approved AI
devices: of the AI-enabled medical devices authorized by the FDA as of 2024,
the majority were approved on the basis of single-site, pre-market studies
with no mandatory external validation against independent datasets and no
continuous post-market performance monitoring. The result is an
\textbf{illusion of safety} — confidence intervals that are statistically
precise but ecologically invalid.

This ecological validity problem is further articulated by
\citet{turchi2025ecological}, whose case study of AI-assisted oral cancer
diagnosis reveals three systematic failures in how AI evaluation is designed:
temporal dynamics of clinical decision-making are not captured by single-point
accuracy metrics; the holistic nature of clinical reasoning (involving context,
prior history, and social cues) is reducible neither to image-based inputs nor
to single-session task performance; and standard performance metrics (accuracy,
AUC) fail to reflect clinical priorities such as false-negative harm asymmetry
or workflow integration burden. These failures are not merely technical — they
reflect a structural misalignment between the epistemology of clinical validation
and the operational realities of clinical environments.

The methodological response to these failures is beginning to emerge.
Prospective, multicenter evaluation designs — exemplified by the TRAC-Pain
protocol \citep{bailes2026trac}, which uses continuous wearable sensor data
combined with ecological momentary assessment over a 12-week longitudinal
cohort — represent a more ecologically grounded approach to AI performance
assessment. However, such designs remain the exception rather than the rule.
Complementary to this, the bias detection literature has advanced subgroup
performance evaluation, adversarial robustness testing, and distribution-shift
benchmarks — such as ShiftCamel and similar frameworks — as essential
components of pre-deployment evaluation. These methods interrogate model
behavior across patient subgroups and under perturbed conditions, revealing
failure modes that aggregate metrics conceal. Yet their adoption remains
inconsistent in regulatory evaluation protocols, and there is no consensus
on mandatory thresholds for fairness or robustness metrics prior to clinical
deployment \citep{pham2025ethical}.

% =============================================================================
% 2.2 GOVERNANCE AND REGULATION — regulatory frameworks and their limits
% =============================================================================

\subsection{Governance and Regulatory Frameworks}
\label{sec:governance}

At the regulatory level, two major frameworks have shaped the governance
landscape for health AI: the FDA's risk-based premarket review pathway — which,
as \citet{warraich2024fda} document, has authorized nearly 1,000 AI-enabled
medical devices — and the European Union's AI Act (2024), which classifies
AI systems in clinical care as high-risk and subjects them to conformity
assessment procedures prior to market entry. Both frameworks share a common
assumption: that governance is a precondition that health systems must
demonstrate before deployment. Against this assumption, the evidence from
\citet{tian2026pilottrap} offers a powerful counterpoint.

Based on an 18-month implementation study of a provincial clinical AI platform
in China, Tian and colleagues argue that governance capacity is not
pre-established but rather \textbf{built through implementation itself}. The
six-module governance framework they propose — covering institutional carrier
formation, infrastructure governance, regulatory and ethical governance,
interdisciplinary coordination, translational scaling, and lifecycle evaluation
— is derived empirically from operational tensions observed during deployment,
not prescribed in advance. This represents a significant departure from the
pre-market authorization model, which treats governance as a gatekeeping
function rather than a developmental process.

The tension between prescriptive regulatory standards and emergent governance
capacity is further highlighted by \citet{abulibdeh2025illusion}, who argue
that the FDA's current framework fails to mandate continuous post-market
performance monitoring, does not require disclosure of training data sources,
and lacks enforceable standards for fairness and bias mitigation. In response,
they call for an \textbf{adaptive, community-engaged regulatory approach} that
integrates mandatory post-market evaluations, algorithmic explainability
requirements, and multi-stakeholder oversight. This is consistent with the
broader ethical governance literature, which emphasizes that AI accountability
cannot rest on technical compliance alone and must embed participatory
mechanisms to ensure that AI systems reflect the values of the communities
they affect \citep{bouderhem2025shaping, abdul2025ai}.

The EU AI Act and the FDA's approach represent two structurally distinct
regulatory postures with different implications for implementation. The EU's
ex-ante model imposes conformity assessment and documentation requirements
before market entry, effectively treating governance as a precondition to be
verified prior to any real-world deployment. Implementation science suggests
this approach faces a fundamental challenge: governance capacity is
context-dependent, organization-specific, and emergent — qualities that cannot
be fully assessed in a pre-market evaluation scenario \citep{tian2026pilottrap}.
The FDA's ex-post posture — authorizing devices on the basis of premarket
studies and relying on post-market surveillance to identify failures — is
likewise problematic because it defers evidence generation to a period when
the device is already in widespread clinical use, creating exposure to the
harms that premarket evaluation was designed to prevent \citep{abulibdeh2025illusion}.
Structured maturity models for AI governance \citep{hussein2026governance} and
layered auditing frameworks linking continuous monitoring to institutional
decision-making \citep{palama2026auditing} represent emerging approaches that
seek to bridge this gap, operationalizing adaptive governance across the
device lifecycle rather than treating it as a gatekeeping checkpoint.

% =============================================================================
% 2.3 IMPLEMENTATION OUTCOMES — from pilot to routine care
% =============================================================================

\subsection{Implementation Outcomes and Barriers}
\label{sec:implementation}

The implementation science literature offers a necessary corrective to the
optimistic narrative that surrounds health AI deployment. The pilot trap
phenomenon documented by \citet{tian2026pilottrap} is one of the most
systematic accounts of why AI pilots fail to scale: the study found that the
governance challenges — unclear accountability boundaries, absent coordination
mechanisms, and inadequate infrastructure governance — are not side effects
of underperforming algorithms but the primary barrier to transition from pilot
to routine clinical care. This finding is corroborated by \citet{keri2026simulated},
who documents in the psychiatric AI context the safety risks that emerge when
AI systems are deployed without adequate validation, monitoring, and human
accountability protocols. In psychiatry specifically, the risks of
\textit{simulated consciousness} — where patients overattribute competence to
AI systems — and of \textit{illusion of attachment} represent a qualitatively
distinct failure mode that standard evaluation frameworks do not yet address.

Against this backdrop, \citet{pham2025ethical} offer a synthesizing perspective
that connects governance to implementation outcomes. Their multi-disciplinary
framework — spanning technologists, healthcare providers, legal experts, and
policymakers — emphasizes that safe and equitable AI deployment requires not
only technical validation but also public engagement, transparent regulatory
processes, and globally harmonized adaptive frameworks. This connects directly
to the participatory governance literature, which holds that the legitimacy
and sustainability of clinical AI ultimately depends on the degree to which
affected communities are meaningfully included in governance processes.

Implementation barriers to health AI integration are multi-dimensional and
operate at organizational, technical, and social levels. At the organizational
level, workforce trust remains a persistent challenge: clinicians who
perceive AI systems as opaque or as threatening their professional autonomy
are less likely to adopt them meaningfully, regardless of technical
performance \citep{keri2026simulated}. Change management deficits — including
inadequate training, misaligned incentive structures (where AI use is not
rewarded or is penalized relative to productivity targets), and insufficient
engagement of frontline staff in deployment decisions — compound this resistance.
At the technical level, data quality inconsistencies, poor interoperability
between AI systems and existing clinical information systems, and inadequate
infrastructure for real-time monitoring create friction that is amplified in
resource-constrained settings. The NASSS framework \citep[Non-Adoption,
Abandonment, Scale-up, Spread, Sustainability]{greenhalgh2017nasps} provides
a structured lens for analyzing these interacting barriers, identifying
disease complexity, organizational readiness, and funding stability as key
determinants of implementation success that are easily overlooked in
technology-centric evaluation designs.

At the social level, algorithmic literacy among both clinicians and patients
mediates how AI outputs are interpreted and acted upon. Low algorithmic
literacy can lead to both over-trust (treating AI recommendations as
authoritative) and under-trust (dismissing useful outputs), each with
distinct clinical consequences. Power dynamics also shape AI adoption:
hierarchical structures in clinical environments can suppress front-line
concerns about AI reliability, while institutional power asymmetries between
technology vendors and healthcare organizations can prevent meaningful
contractual accountability for AI failures \citep{abulibdeh2025illusion}.
Addressing these social barriers requires governance mechanisms that are
not merely procedural but are embedded in organizational cultures that
value critical engagement with AI systems and maintain human accountability
as a non-negotiable condition of deployment \citep{bouderhem2025shaping,
abdul2025ai, hussein2026governance}.